\documentclass[a4paper, 12pt]{article}
% math symbols
\usepackage{amssymb}
\usepackage{amsmath}
\usepackage{mathrsfs}
\usepackage{mathseries}


\usepackage[margin = 2cm]{geometry}

\tolerance = 1000
\emergencystretch = 0.74cm



\pagestyle{empty}
\parindent = 0mm

\renewcommand{\coursetitle}{DM/ML}
\setcounter{curtask}{1}

\setmathstyle{\includegraphics[scale = 0.05]{../pics/utia-rest.png}}{Мат. логика}{3 курс}


\begin{document}

\task{
    Покажите, что в интерпретации $(\mathbb{R}, 0, 1, +,\, \cdot \, , <, =)$ предикат $x = \alpha$, где
    $\alpha$~--- вещественное число, выразим тогда, и только тогда, когда $\alpha$~--- алгебраическое
    число.
}


\libdefinition{real-closed-field}

\task{
    Выведите в теории линейно упорядоченных полей:
    $$
        x < y \longrightarrow \exists u\, x < u < y.
    $$
}


\begin{proposition}
    Для любого линейно упорядоченного поля $\mathbb{F}$ следующие условия эквивалентны:
    \begin{enumcyr}
        \item $\mathbb{F}$ вещественно замкнуто;
        \item $\mathbb{F}$ как поле не является алгебраически замкнутым, но его расширение
            $\mathbb{F}[\sqrt{-1}]$ алгебраически замкнуто.
    \end{enumcyr}
\end{proposition}

\task{
    Покажите, что для любого линейно упорядоченного поля $\mathbb{F}$ следующие условия эквивалентны:
    \begin{enumcyr}
        \item $\mathbb{F}$ вещественно замкнуто;
        \item в $\mathbb{F}$ истинно $\forall x\, (0 < x \rightarrow \exists u\, x = u \cdot u)$, а
            также все $\sigma$-формулы вида:
            $$
                \forall \overline{u}\, \forall x\, \forall y\,
                (\mathrm{p}_n(x; \overline{u}) \cdot \mathrm{p}_n(y; \overline{u}) < 0 \rightarrow
                \exists z\, (x < z < y \wedge \mathrm{p}_n(z; \overline{u}) = 0)),
            $$
            т.е.\ для $\mathbb{F}$ верна <<полиномиальная версия>> теоремы о среднем значении.
    \end{enumcyr}
}

\libproblem{math-logic}{rcf-nonstandard-point}
\libproblem{math-logic}{acf-char-0-big-char}

\libproblem{math-logic}{red-blue-dlo-count}
\libproblem{math-logic}{predicate-strategy-exists}




\end{document}



%%% Local Variables:
%%% mode: latex
%%% TeX-master: t
%%% End:
